\documentclass[11pt,a4paper,oneside,notitlepage]{book}
\usepackage[english,dutch]{babel}

% marges aanpassen
% (opmerking: moet *voor* inclusie van fancyhdr package komen)
\setlength{\hoffset}{-1in}
\setlength{\voffset}{-1in}
\setlength{\topmargin}{2cm}
\setlength{\headheight}{0.5cm}
\setlength{\headsep}{1cm}
\setlength{\oddsidemargin}{3.5cm}
\setlength{\evensidemargin}{3.5cm}
\setlength{\textwidth}{16cm}
\setlength{\textheight}{23.3cm}
\setlength{\footskip}{1.5cm}

\usepackage[a4paper,plainpages=false]{hyperref}    % Om hyperlinks te hebben in het pdfdocument.
\usepackage{url}
\usepackage{fancyhdr}
\usepackage{graphicx}
%custom packages
\usepackage{booktabs}
\usepackage{longtable}
\usepackage{ltcaption}
\usepackage{listings}
\lstset{
	language=C,
	basicstyle=\small\sffamily,
	breaklines=true,
	numberstyle=\tiny,
	frame=tb,
	columns=fullflexible,
	showstringspaces=false
}
\newcommand{\quotes}[1]{``#1''}
% \usepackage[colorlinks]{hyperref}

\pagestyle{fancy}

\renewcommand{\chaptermark}[1]{\markright{\MakeUppercase{#1}}}
\renewcommand{\sectionmark}[1]{\markright{\thesection~#1}}

\newcommand{\headerfmt}[1]{\textsl{\textsf{#1}}}
\newcommand{\headerfmtpage}[1]{\textsf{#1}}

\fancyhf{}
\fancyhead[LE,RO]{\headerfmtpage{\thepage}}
\fancyhead[LO]{\headerfmt{\rightmark}}
\fancyhead[RE]{\headerfmt{\leftmark}}
\renewcommand{\headrulewidth}{0.5pt}
\renewcommand{\footrulewidth}{0pt}


\fancypagestyle{plain}{ % eerste bladzijde van een hoofdstuk
  \fancyhf{}
  \fancyhead[LE,RO]{\headerfmtpage{\thepage}}
  \fancyhead[LO]{\headerfmt{\rightmark}}
  \fancyhead[RE]{\headerfmt{\leftmark}}
  \renewcommand{\headrulewidth}{0.5pt}
  \renewcommand{\footrulewidth}{0pt}
}

% anderhalve interlinie (opm: titelblad gaat uit van 1.5)
\renewcommand{\baselinestretch}{1.5}

% indien LaTeX niet goed splitst, neem je woord hierin op, of evt om splitsen 
% te voorkomen
\hyphenation{ditmagnooitgesplitstworden dit-woord-splitst-hier}

\begin{document}

% titelblad (voor kaft)
%  Titelblad

% Opmerking: gaat uit van een \baselinestretch waarde van 1.5 (die moet
% ingesteld worden voor het begin van de document environment)

\begin{titlepage}

\setlength{\hoffset}{-1in}
\setlength{\voffset}{-1in}
\setlength{\topmargin}{1.5cm}
\setlength{\headheight}{0.5cm}
\setlength{\headsep}{1cm}
\setlength{\oddsidemargin}{3cm}
\setlength{\evensidemargin}{3cm}
\setlength{\footskip}{1.5cm}
\enlargethispage{1cm}
% \textwidth en \textheight hier aanpassen blijkt niet te werken

\fontsize{12pt}{14pt}
\selectfont

\begin{center}

\includegraphics[height=2cm]{fig/ruglogo}

\vspace{0.5cm}

Universiteit Gent\\
Faculteit Ingenieurswetenschappen en Architectuur\\

\vspace{4.5cm}

\fontseries{bx}
\fontsize{17.28pt}{21pt}
\selectfont

Verbeterde aanbevelingssystemen 
\\op basis van content herkenning in tekst


\fontseries{m}
\fontsize{12pt}{14pt}
\selectfont

\vspace{.6cm}

door 

\vspace{.4cm}

Bjorn VANDENBUSSCHE

\vspace{3.5cm}

Promotor: prof.~dr.~ir.~Luc~Martens\\
Begeleider: dr.~ir.~Toon~De Pessemier\\

\vspace{2cm}

Scriptie ingediend tot het behalen van de academische graad van\\ 
Master of Science in de industri\"ele wetenschappen: informatica

\vspace{1cm}

Academiejaar 2015--2016

\end{center}
\end{titlepage}


% lege pagina (!!)

% titelblad (!!)

% geen paginanummering tot we aan de inhoudsopgave komen
\pagestyle{empty}

% voorwoord met dankwoord en toelating tot bruikleen (ondertekend)
%  Voorwoord (dankwoord) en toelating tot bruikleen

\newpage

\noindent \textbf{\huge Voorwoord}

\vspace{1.5cm}

\noindent
Hier komt het voorwoord -- Dit is de voorlopige versie zoals ingediend op 30/3/2016

\addvspace{4cm}

\noindent Bjorn Vandenbussche, juni 2016\newpage

\noindent \textbf{\huge Toelating tot bruikleen}

\vspace{1.5cm}

\noindent
``De auteur geeft de toelating deze scriptie voor consultatie beschikbaar
te stellen en delen van de scriptie te kopi\"eren voor persoonlijk
gebruik.\\
Elk ander gebruik valt onder de beperkingen van het auteursrecht,
in het bijzonder met betrekking tot de verplichting de bron uitdrukkelijk
te vermelden bij het aanhalen van resultaten uit deze scriptie.''

\addvspace{4cm}

\noindent Bjorn Vandenbussche, juni 2016


% overzicht
%  Overzichtsbladzijde met samenvatting

\newpage

{
\setlength{\baselineskip}{14pt}
\setlength{\parindent}{0pt}
\setlength{\parskip}{8pt}

\begin{center}

\noindent \textbf{\huge
Verbeterde aanbevelingssystemen\\[8pt]op basis van content recognition in tekst
}

door 

Bjorn VANDENBUSSCHE

Scriptie ingediend tot het behalen van de academische graad van\\ 
Master of Science in de industri\"ele wetenschappen: informatica

Academiejaar 2015--2016

Promotor: Prof.~L.~MARTENS\\
Scriptiebegeleider: Dr.~ir.~T.~DE PESSEMIER\\

Faculteit Ingenieurswetenschappen en Architectuur\\
Universiteit Gent

Vakgroep Informatietechnologie

\end{center}

\section*{Samenvatting}

De klassieke aanbevelingssystemen werken op basis van consumptiegedrag (beoordelingen, aankopen, klikgedrag), soms aangevuld met informatie over de content (metadata) zoals sleutelwoorden, auteur, categorie\"en, samenvatting,... Het probleem bij deze systemen is dat de beschrijving niet altijd perfect is.  Aanbevelingssystemen zullen in dat geval niet goed werken, gezien deze een exacte overeenkomst van de metadata velden verwachten of op zijn minst een overlap in gebruikte woorden.

Het doel van dit onderzoek is om te achterhalen welke vormen van content recognition het aanbevelingssysteem kunnen helpen om nuttige aanbevelingen te leveren. Verschillende informatie-extractietechnieken kunnen ervoor zorgen dat het aanbevelingssysteem een beter beeld krijgt van de voorkeur van de gebruiker door werkelijk rekening te houden met de inhoudelijke aspecten van de items die aanbevolen worden.

Daartoe zijn verschillende vormen van content recognition uitgetest, voornamelijk rond classificatie en clustering. Door verschillende items binnen dezelfde categorie te classificeren kan de context al heel wat duidelijker afgeleid worden. Om vervolgens bepaalde \textit{topics} binnen categorie\"en of samenhorende items over verschillende categorie\"en heen te vinden wordt clustering toegepast. Door zowel de informatie van de classificatie als van de clustering te combineren krijgt het aanbevelingssysteem een beter beeld van de inhoud van de items en worden bijgevolg betere aanbevelingen gegenereerd.

\section*{Trefwoorden}

text mining, big data, aanbevelingssystemen, machine learning

}

\newpage % strikt noodzakelijk om een header op deze blz. te vermijden

%  Engelse variant


{
\setlength{\baselineskip}{14pt}
\setlength{\parindent}{0pt}
\setlength{\parskip}{8pt}

\begin{center}

\noindent \textbf{\huge
Better recommender systems\\[8pt]based on content recognition in text
}

by 

Bjorn VANDENBUSSCHE

Thesis submitted to obtain the academic degree of\\
Master of Science in Industrial Sciences: Informatics

Academic year 2015--2016

Supervisor: Prof.~L.~MARTENS\\
Mentor: Dr.~ir.~T.~DE PESSEMIER\\

Faculty of Engineering and Architecture\\
University of Ghent

Department of Information Technology

\end{center}

\section*{Abstract}

Classical recommender systems work based on consumption behavior (ratings, purchases, clicks), sometimes supplemented by information about content (metadata) lika keywords, author, category, resume, etc. The problem with these systems is that the description isn't always perfect. Recommender systems won't work well in those cases, as they expect an exact match of metadata or at least an overlap in used words.

The goal of this research is to find out what kinds of content recognition can help the recommender system to deliver useful recommendations. Different techniques for information extraction can make sure the recommender gets a better view on the preferences of a user, by really taking the content of items being recommended into account.

Therefore multiple forms of content recognition where tested, mostly concerning classification and clustering. Classifying multiple items in the same category makes it possible to deduct context of an item. To find different topics inside of a category or to find related items over the boundaries of the known categories, clustering is applied. By combining the information from both classification and clustering the recommender systems gets a better view on the content of the items and is thereby capable of generating better recommendations.

\section*{Keywords}

text mining, big data, recommender systems, machine learning

}

\newpage % strikt noodzakelijk om een header op deze blz. te vermijden


\pagestyle{fancy}
\frontmatter

% inhoudstafel
\tableofcontents

% opmaak voor het eigenlijke boek; onderstaande lijnen
% weglaten als de eerste regel van een nieuwe alinea moet
% inspringen in plaats van extra tussenruimte
\setlength{\parindent}{0pt}
\setlength{\parskip}{0.5\baselineskip plus 0.5ex minus 0.2ex}
\setlength{\parskip}{1ex plus 0.5ex minus 0.2ex}

% hoofdstukken
\mainmatter

% hier worden de hoofdstukken ingevoegd (\includes)
\chapter{Inleiding}

\section{Bestaande situatie}
De enorme hoeveelheid informatie op het internet confronteert de eindgebruiker met het probleem van overaanbod.
Hoewel alle content beschikbaar is, is het voor de gebruiker moeilijk om de meest geschikte, de meest interessante
content terug te vinden. 

De klassieke sleutelwoord gebaseerde zoeksystemen bieden hierbij slechts een gedeeltelijke
oplossing. Voorkeuren van de gebruiker, eerdere ervaringen, en reeds geconsumeerde content worden niet in
rekening gebracht bij de klassieke zoeksystemen. Daarom zijn aanbevelingssystemen een nuttig hulpmiddel voor het
ontdekken, selecteren en consumeren van content. Vele online diensten zoals Amazon \cite{Everything2012} en Netflix \cite{Bennett2007} hebben
reeds een aanbevelingssysteem dat zijn nut al meermaals bewezen heeft. 

De klassieke aanbevelingssystemen werken op basis van consumptiegedrag (beoordelingen, aankopen, klikgedrag),
soms aangevuld met informatie over de content (metadata) zoals sleutelwoorden, auteur, categorie\"en, samenvatting,... Het probleem bij deze systemen is dat de beschrijving niet altijd perfect is. Zo kunnen er schrijffouten optreden, synoniemen of hyponiemen voorkomen, of een ander lexicon of thesaurus gebruikt zijn.

Aanbevelingssystemen zullen in dat geval niet goed werken, gezien deze een exacte overeenkomst van de metadata
velden verwachten of op zijn minst een overlap in gebruikte woorden. Zo zullen bij klassieke aanbevelingssystemen
twee synoniemen ten onrechte als twee verschillende concepten beschouwd worden. Bepaalde verbanden tussen
gerelateerde content items zullen dus nooit door een klassiek aanbevelingssysteem ontdekt worden.
Op het web is er echter meer informatie beschikbaar dan typisch gebruikt wordt in aanbevelingssystemen. Door het
“begrijpen” van de inhoud, kunnen aanbevelingssystemen meer nauwkeurige suggesties doen.

\section{Doelstelling}
Het doel van deze masterproef is de ontwikkeling van een aanbevelingssysteem die gebaseerd is op content
recognition in tekst. Via content recognition technieken kan informatie uit teksten gehaald worden die als input voor
een aanbevelingssysteem gebruikt kunnen worden. Informatie-extractie kan ervoor zorgen dat het
aanbevelingssysteem een beter beeld krijgt van de voorkeuren van de gebruiker door werkelijk rekening te houden
met de inhoudelijke aspecten van de items die aanbevolen worden. Tevens zal content recognition ervoor zorgen dat
de afhankelijkheid van metadata in aanbevelingssystemen minder sterk is, of zelfs verdwijnt.

\section{Overzicht}
TODO: hs1 -> ..., hs2 -> ...
\chapter{Achtergrond}

De opgeslagen informatie neemt elke dag enorm toe. Ontdekken van patronen en trends uit deze steeds groeiende data is een taak die steeds moeilijker wordt. Bepaalde technieken kunnen gebruikt worden om het probleem op te lossen. De basistechniek daarbij is \textit{data mining} (zie \ref{data-mining}), waarvan toepassingen o.a. \textit{text mining} (zie \ref{text-mining}) en \textit{web mining} zijn \cite{Nasa2012}. 

Dit hoofdstuk voorziet de nodige achtergrondinformatie om de werking van dergelijke text mining systemen te begrijpen. Eerst wordt data mining als algemeen, abstract begrip omschreven. Uiteraard bestaan er genoeg goede publicaties die verder op dit onderwerp ingaan. Vervolgens komt zijn afgeleide toepassing, text mining, aan bod. De belangrijkste topics binnen dit domein worden ten slotte nader uitgelicht. 

\section{Data mining}\label{data-mining}
Data mining werd in 1980 voor het eerst gebruikt om data om te zetten naar kennis. Het doel van data mining is om impliciete, vooraf onbekende trends en patronen uit databases te halen. Daarbij wordt gebruik gemaakt van verschillende technieken: classificatie, clustering, neurale netwerken, beslissingsbomen, etc. 
\\Data mining is deel van het \textit{knowledge discovering process} of \textit{knowledge discovery in databases} (KDD). Dit proces bestaat uit verschillende stappen \cite{Fayyad1996}:
\begin{enumerate}
	\item Begrijpen van business: de objectieven en verwachtingen worden gedefinieerd.
	\item Begrijpen van data: data wordt uit het data warehouse geselecteerd op basis van de gedefinieerde objectieven en verwachtingen.
	\item Voorbereiden van data: de kwaliteit van de data wordt verbeterd.
	\item Modelleren van data: een data mining algoritme wordt geselecteerd en toegepast op de data die voorbereid werd in de vorige stap.
	\item Evaluatie: de relaties en patronen worden geanalyseerd en geldige patronen volgens de vooraf opgestelde doelen worden geselecteerd.
	\item Visuele representatie: de kennis die ontdekt is wordt visueel voorgesteld. Deze resultaten kunnen opgeslagen en samengevoegd worden, om de business vooruit te helpen.
\end{enumerate}

\section{Text mining}\label{text-mining}
\textit{Content recognition} of \textit{text mining} verwijst naar de extractie van interessante informatie en kennis uit ongestructureerde tekst. Het probeert om de verborgen informatie te onthullen door middel van methodes die enerzijds met een groot aantal woorden en structuren in natuurlijke taal kunnen omgaan en anderzijds een zekere vaagheid en onzekerheid kunnen verwerken. Text mining kan naast met gestructureerde data (zoals de data die in data mining verwerkt wordt) ook werken met ongestructureerde of semi-gestructureerde data zoals e-mails, volledige tekstdocumenten, HTML bestanden, etc. Text mining wordt daarom beschreven als een interdisciplinaire methode op basis van \textit{information retrieval}\footnote{Het vinden van documenten die antwoorden bevatten op vragen, maar niet het vinden van de antwoorden zelf.},\textit{ machine learning}, statistiek, computationele taalkunde en vooral data mining  \cite{Hotho2005}. 

Om grote collecties van documenten te verwerken, moeten tekstdocumenten vooraf verwerkt worden om de informatie op te slaan in een datastructuur die beter geschikt is dan een tekstbestand. De meeste text mining-methodes gaan ervan uit dat een tekstdocument voorgesteld kan worden als een set van woorden (\textit{bag-of-words} representatie, zie \ref{bag-of-words}). Ondertussen bestaan echter methodes die proberen om de syntactische structuur of de semantiek in de tekst uit te buiten. Op die manier kan de belangrijkheid van een woord achterhaald worden. Hiervoor wordt vaak een vector representatie gebruikt, die voor elk woord een numeriek gewicht bijhoudt. Enkele belangrijke modellen die hiervoor gebruikt worden zijn het vector space model \cite{Salton1975}, het probabilistische model \cite{ROBERTSON1977} en het logische model \cite{Rigsbergen1986}. Het vector space model wordt verder besproken in \ref{vector-space-model}.

\section{Data mining methodes voor tekst}
De hoofdreden om data mining methodes in te zetten voor documentencollecties is om ze te structureren. Een structuur kan het gemakkelijker maken voor de gebruiker om een documentencollectie te raadplegen. Bestaande methodes die toestaan om documentencollecties te structureren proberen om kernwoorden aan documenten te koppelen op basis van een gegeven set kernwoorden (via classificatie of categorisatie, zie \ref{classificatie}) of proberen de documentencollectie automatisch te structureren in groepen van gelijkaardige documenten (\textit{clustering}, zie \ref{clustering}).


\subsection{Classificatie}\label{classificatie}
Classificatie of categorisatie heeft als doel om voorgedefinieerde klasses toe te kennen aan tekstdocumenten. Het is een \textit{supervised} techniek\footnote{Een set van input-output voorbeelden worden gebruikt om het model te trainen, om zo nieuwe documenten te kunnen classifiseren.} die de \textit{classifier} traint op basis van gekende voorbeelden en zijn opgesteld model vervolgens gebruikt om ongekende voorbeelden automatisch te categoriseren \cite{Nasa2012}. 



\subsection{Clustering}\label{clustering}
Clustering is een \textit{unsupervised} techniek waar geen patronen voorgedefinieerd zijn. De methode is gebasseerd op een concept dat gelijkaardige documenten of teksten in dezelfde cluster groepeert. Elke cluster bevat dus een aantal documenten. De clustering wordt als beter beschouwd indien de inhoud van de documenten intra-cluster meer gelijkheid vertoont dan de inhoud van de documenten inter-cluster.

Clustering wordt gebruikt om gelijkaardige documenten te groeperen, maar verschilt van classificatie omdat documenten bij clustering on-the-fly ingedeeld worden in clusters, in plaats van in vooraf gedefinieerde klasses of topics.

\subsection{Aanbevelingssyteem}
%TODO uitbreiden
Er bestaan verschillende strategie\"en om aanbevelingen te genereren. Deze strategie\"en worden onderverdeelt in volgende categorie\"en \cite{Adomavicius2005}: 
\paragraph{Content-based aanbevelingen}
De gebruiker krijgt aanbevelingen gelijkaardig aan de items die hij vroeger prefereerde.
\paragraph{Collaborative filtering}
De gebruiker krijgt aanbevelingen die andere gebruikers met gelijkaardige smaak en voorkeur goed vonden.
\paragraph{Hybride systemen}
Deze methodes combineren content-based aanbevelingen en collaborative filtering technieken.
\chapter{Architectuur van het ontworpen systeem}

\section{Algemeen overzicht}

Alle onderdelen van het ontworpen systeem zijn onderling onafhankelijk van elkaar. Dit maakt een robuuste, gedistribueerde eenheid waaraan relatief eenvoudig nieuwe componenten kunnen worden toegevoegd. De flow van het programma wordt gecontroleerd door de \textit{FlowController}. Deze zal de flow van het programma bijhouden en alle interactie met de gebruiker verzorgen. Met over de flow en de controller volgt in hoofdstuk \ref{flowcontroller}.

\textbf{TODO: Wat mogen we verwachten in dit hoofdstuk? Geef overzicht}
\subsection{Metadata extractie met Apache Tika}
Om met verschillende documentencollecties te kunnen werken binnen het ontworpen systeem, wordt gebruik gemaakt van Apache Tika \cite{tika}. Tika laat toe om metadata van een tekst af te leiden. Die metadata kan taal of formaat van het document in kwestie zijn, welke ons een beter beeld geven van o.a. het stemmingsalgoritme dat gebruikt moet worden of laat toe om verschillende soorten documentformaten (pdf, doc, docx, csv, html, etc) om te vormen tot platte tekst zonder franjes. Deze tekst kan veel beter ge\"interpreteerd worden door WEKA en zal dus betere resultaten opleveren voor o.a. classificatie.

Apache tika wordt aangeboden onder de vorm van een REST-API, gebouwd op het Jersey-framework in Java. Het Jersey-framework is gebaseerd op JAX-RS en laat toe om verschillende types data aan te bieden aan de gebruiker. Het wordt gedeployed naar een Glassfish-omgeving, die vervolgens kan worden aangesproken door gebruikers van de API.

Op zich bevat de API slechts twee verschillende \textit{endpoints}; \quotes{/test} controleert of de API wel degelijk reageert op calls, terwijl \quotes{/upload} een document aanvaardt in de vorm van een tekst en die gaat parsen naar platte tekst. Het parsen bestaat uit twee delen. Het eerste deel is de taaldetectie, welke door de standaard \textit{LanguageIdentifier} van het Tika framework wordt uitgevoerd. Het tweede deel vervolgens is het effectief parsen van de tekst. Hiervoor wordt de \textit{AutoDetectParser} van Tika gebruikt. Deze kiest voor elk document de parser die voor dat type het beste resultaat zal opleveren.
De resultaten van de parser worden teruggestuurd naar de initiator van de connectie in de vorm van JSON:

\begin{lstlisting}
{
"language":"nl",
"title":"5-000-deelnemers-opname-clip-klimaatsverandering\n",
"body":"Op het Klein Strand in Oostende zijn vanmiddag naar schatting 5.000 mensen samengekomen op een klimaathappening die onderdeel was van de internationale campagne 'The Big Ask'. De deelnemers figureerden in een filmpje van regisseur Nic Baltazar en Friends of the Earth vzw.Filmpje YoutubeOp het Klein Strand vormden de deelnemers letters en slogans die vanuit de hoogte werden opgenomen en later zullen te zien zijn in een filmpje op Youtube. De actie 'Sos Klimaat : bouw een dijk tegen klimaatsverandering' kreeg de steun van talrijke BV's die als dj de massa animeerden. Onder meer Zohra, Flip Kowlier, Gabriel Rios en Adriaan Van den Hoof verleenden op het podium hun medewerking.Politici wakkerschuddenInitiatiefnemer Nic Baltazar was in elk geval tevreden met de opkomst en de opnames en verwacht dat de clip op Youtube voor een half miljoen hits zal zorgen. Hij hoopt dat de actie en de clip politici zal wakker schudden omtrent de problematiek. (belga/ep)\n"
}
\end{lstlisting} 


\subsection{Preprocessing van de documentencollectie}

\subsection{Classificatie}

\subsection{Clustering}

\subsection{Recommender system}

\chapter{Overzicht werking}

\section{Algemeen}
Gebruik Java/Perl/... uitleggen

\section{WEKA}
WEKA (Waikato Environment for Knowledge Analysis) is een tool voor datamining ontwikkeld in de programmeertaal Java. Het bestaat uit een grafische werkomgeving  voor het uitvoeren van de nodige stappen bij datamining, een CLI en is volledig ondersteund voor gebruik binnen een eigen Java-programma. Het helpt bij de preprocessing van data en bevat heel wat algoritmes voor clustering, classificatie, regressie-analyse, visualisatie en feature selection.

De versie van WEKA die gebruikt wordt in deze masterproef is 3.6.13, de op heden laatste stabiele versie.

\textbf{TODO: meer + referentie}

\section{Metadata extractie met Apache Tika}
Om met verschillende documentencollecties te kunnen werken binnen het ontworpen systeem, wordt gebruik gemaakt van Apache Tika \cite{tika}. Tika laat toe om metadata van een tekst af te leiden. Die metadata kan taal of formaat van het document in kwestie zijn, welke ons een beter beeld geven van o.a. het stemmingsalgoritme dat gebruikt moet worden of laat toe om verschillende soorten documentformaten (pdf, doc, docx, csv, html, etc) om te vormen tot platte tekst zonder franjes. Deze tekst kan veel beter ge\"interpreteerd worden door WEKA en zal dus betere resultaten opleveren voor o.a. classificatie.

Apache tika wordt aangeboden onder de vorm van een REST-API, gebouwd op het Jersey-framework in Java. Het Jersey-framework is gebaseerd op JAX-RS en laat toe om verschillende types data aan te bieden aan de gebruiker. Het wordt gedeployed naar een Glassfish-omgeving, die vervolgens kan worden aangesproken door gebruikers van de API.

Op zich bevat de API slechts twee verschillende \textit{endpoints}; \quotes{/test} controleert of de API wel degelijk reageert op calls, terwijl \quotes{/upload} een document aanvaardt in de vorm van een tekst en die gaat parsen naar platte tekst. Het parsen bestaat uit twee delen. Het eerste deel is de taaldetectie, welke door de standaard \textit{LanguageIdentifier} van het Tika framework wordt uitgevoerd. Het tweede deel vervolgens is het effectief parsen van de tekst. Hiervoor wordt de \textit{AutoDetectParser} van Tika gebruikt. Deze kiest voor elk document de parser die voor dat type het beste resultaat zal opleveren.
De resultaten van de parser worden teruggestuurd naar de initiator van de connectie in de vorm van JSON:

\begin{lstlisting}
{
"language":"nl",
"title":"5-000-deelnemers-opname-clip-klimaatsverandering\n",
"body":"Op het Klein Strand in Oostende zijn vanmiddag naar schatting 5.000 mensen samengekomen op een klimaathappening die onderdeel was van de internationale campagne 'The Big Ask'. De deelnemers figureerden in een filmpje van regisseur Nic Baltazar en Friends of the Earth vzw.Filmpje YoutubeOp het Klein Strand vormden de deelnemers letters en slogans die vanuit de hoogte werden opgenomen en later zullen te zien zijn in een filmpje op Youtube. De actie 'Sos Klimaat : bouw een dijk tegen klimaatsverandering' kreeg de steun van talrijke BV's die als dj de massa animeerden. Onder meer Zohra, Flip Kowlier, Gabriel Rios en Adriaan Van den Hoof verleenden op het podium hun medewerking.Politici wakkerschuddenInitiatiefnemer Nic Baltazar was in elk geval tevreden met de opkomst en de opnames en verwacht dat de clip op Youtube voor een half miljoen hits zal zorgen. Hij hoopt dat de actie en de clip politici zal wakker schudden omtrent de problematiek. (belga/ep)\n"
}
\end{lstlisting} 

\section{Preprocessing van de documentencollectie}\label{preprocessing}
De datasets van Het Laatste Nieuws bevatten, afhankelijk van de tijdsperiode, 11 tot 12 hoofdcategorie\"en. Voor de dataset van januari 2008 zijn deze hoofdcategorie\"en met hun aandeel in het totale aantal artikels te vinden in tabel \ref{tab:hln-2008-01-cat}. Elke hoofdcategorie bevat bovendien een aantal subcategorie\"en. Elk artikel behoort bijgevolg toe aan twee categorie\"en. 

De eerste stap, voor we de classificatie en clustering toepassen, omvat het preprocessen van alle documenten uit de documentencollectie. In WEKA wordt hiervoor de functie \textit{string-to-wordvector} gebruikt. Deze functie converteert de teksten naar een set van attributen die het aandeel van dat attribuut voorstellen in de documentencollectie (zie \ref{vector-space-model}). 

De string-to-word functie in WEKA neemt een aantal parameters. Een eerste parameter is de tokenizer functie. Deze bepaalt hoe een tekst wordt opgesplitst in woorden. Hier wordt gebruik gemaakt van een N-Gram tokenizer. Deze geeft de mogelijkheid om eventueel twee of drie woorden samen te nemen en als \'e\'en woord in de documentencollectie te zien.

Ten tweede kan gekozen worden voor term frequency-inverse document frequency (tf-idf), al dan niet gecombineerd met een normalisatie. De achterliggende logica zit vervat in formule \ref{eq:tfidf}. Op deze manier worden termen die belangrijk zijn voor een document, maar minder vaak of niet in de volledige documentencollectie voorkomen, een groter gewicht toegekend in de vector representatie. 

Een derde parameter die ingesteld kan worden is de het gebruikte stemmingsalgoritme. Hier werd een Nederlandstalige stemmer gebruikt, zoals beschreven in \cite{Kraaij1994}. Deze stemmer is een uitbreiding van het originele stemmingsalgoritme van Porter, zoals voorgesteld in \cite{Porter1980}, naar de Nederlandse taal. 

Een vierde parameter is de stopwoordfilter. Hiervoor wordt een lijst met stopwoorden gebruikt voor de Nederlandse taal. De stopwoordfilter zorgt ervoor dat woorden die veel voorkomen in een taal, maar relatief weinig inhoud hebben, niet overwogen worden om op te nemen in het woordenboek. Denk daarbij aan 'de', 'het', 'daar', 'dan', etc. Na het opsplitsen van de zin naar woorden (tokenization, zie \ref{tokenization}) kunnen deze woorden weggefilterd worden, zodat ze geen negatieve invloed hebben op het uiteindelijk resultaat.

Een vijfde parameter laat toe om de frequentie van een woord uit te drukken, in plaats van binaire aanwezigheid. Hierdoor wordt de belangrijkheid van bepaalde woorden soms meer in de verf gezet. Wel is het opletten dat sommige minder belangrijke woorden zo niet een grotere rol krijgen dan ze eigenlijk verdienen.

De laatste parameter is het aantal woorden dat behouden wordt per klasse uit de documentencollectie. Dit is natuurlijk een belangrijke parameter, omdat meer betekenisvolle woorden in het woordenboek naar een betere classificatie en/of clustering zullen leiden. Als we echter te veel woorden opnemen, waardoor er een bepaalde ruis gaat optreden, krijgen we ook slechts een suboptimale werking. Het is echter moeilijk om zomaar te bepalen hoeveel woorden een optimale classificatie zullen geven bij elke methode. Bovendien is het mogelijk dat het voor onze dataset niet nuttig is om een stemmer of een stopwoordfilter te gaan gebruiken. Daarom wordt in \ref{result-clas} een benchmark opgesteld die het optimale aantal woorden gaat onderzoeken bij classificatie. Verder wordt ook het  nut van de stemmer en van de stopwoordfilter verder onderzocht. 

\textbf{TODO: hoe doen we dit bij clustering? zelfde?\\-- wrs niet, wel nood aan tf-idf, normalisatie, word count (freq)}

\section{Classificatie}\label{result-clas}
Een eerste techniek die toegepast wordt om een aanbevelingssysteem een mogelijk beter beeld te geven van de documentencollectie, is classificatie. Door een artikel automatisch aan een categorie te kunnen toewijzen, krijgen we al een veel duidelijk beeld van waarover dit artikel nu precies handelt. Als we zowel de hoofd- als subcategorie van een artikel kunnen bepalen, hebben we al een eerste sterke indicatie of het artikel al dan niet interessant kan zijn voor de gebruiker en dus aanbevolen zal worden door het aanbevelingssysteem. 

WEKA ondersteunt echter slechts \'e\'en niveau van classificatie. Om dit probleem op te lossen gaan we eerst de artikels indelen in de hoofdcategorie\"en en pas daarna in de subcategorie\"en. Omdat het mogelijk is dat voor deze gevallen een andere classifier en een ander aantal woorden optimaal is, gaan we deze apart testen. De preprocessing blijft voor alle gevallen gelijk. Wel gaan we eerst uitzoeken welke parameters van de preprocessing het beste beeld geven voor enkele na\"ieve classifiers. 

\subsection{Parameters voor preprocessing van data}

Om de optimale parameters voor de preprocessing te bepalen wordt een classificatie uitgevoerd met de verwerkte dataset. We vari\"eren de parameters en bepalen zo de optimale instellingen. De dataset die gebruikt wordt is deze van Het Laatste Nieuws, januari 2008. De classifiers die gebruikt worden om de bekomen woordvector te testen zijn "Naive Bayes$"$ en "Naive Bayes Multinomial". 

Om een betrouwbaar resultaat te bekomen worden de tests met een \textit{10-fold cross validation}\label{10-fold-cv} uitgevoerd. Dit is een techniek om voorspellende modellen te evalueren door de originele dataset te verdelen in een trainingsset om het model te trainen en een testset om deze te valideren. Bij 10-fold cross-validation\label{cross-val} wordt de originele dataset random verdeeld in 10 subsets. Daarvan wordt \'e\'en subset behouden ter validatie van het model. De overige 9 subsets worden gebruikt als trainingsdata. Het cross-validation proces wordt 10 keer herhaald, waarbij elk van de subsets exact \'e\'en keer gebruikt worden als validatieset. De 10 resultaten worden gecombineerd (uitgemiddeld) om een enkele voorspelling te maken. Het grote voordeel van deze methode is dat alle observaties zowel gebruikt worden als trainings- en validatiedata. 

Om de resultaten onafhankelijk te stellen van het aantal woorden dat behouden wordt van de documentencollectie, gebruiken we de standaard 1000 woorden die WEKA voorstelt te behouden. Verder worden de verschillende parameters al dan niet ingesteld, zoals te zien in tabel \ref{tab:preprocessing-params}. Hierbij staat "1" voor gebruikt en "0" voor niet gebruikt. 


\newpage
\begin{longtable}{rrrrrrr}
	\caption{Gebruik van verschillende parameters bij preprocessing van de documentencollectie en hun effect op het percentage correct geclassificeerde artikels volgens de "Naive Bayes$"$ (NB) en "Naive Bayes Multinomial" (NBM) classifiers.}\\
	\toprule
	\multicolumn{1}{c}{\textbf{tf-idf}} & \multicolumn{1}{c}{\textbf{Stemming}} & \multicolumn{1}{c}{\textbf{Stopwoorden}} & \multicolumn{1}{c}{\textbf{Frequentie}} & \multicolumn{1}{c}{\textbf{Tokenizer}} & \multicolumn{1}{c}{\textbf{NB}} & \multicolumn{1}{c}{\textbf{NBM}} \\
	\midrule
	0     & 0     & 0     & 0     & 1     & 69.85\% & 74.39\% \\
	0     & 0     & 0     & 0     & 2     & 67.62\% & 72.34\% \\
	0     & 0     & 0     & 0     & 3     & 67.44\% & 72.15\% \\
	0     & 0     & 0     & 1     & 1     & 56.06\% & 66.39\% \\
	0     & 0     & 0     & 1     & 2     & 54.26\% & 64.35\% \\
	0     & 0     & 0     & 1     & 3     & 54.24\% & 64.30\% \\
	0     & 0     & 1     & 0     & 1     & 71.50\% & 75.71\% \\
	0     & 0     & 1     & 0     & 2     & 69.26\% & 73.41\% \\
	0     & 0     & 1     & 0     & 3     & 68.91\% & 73.09\% \\
	0     & 0     & 1     & 1     & 1     & 58.01\% & 68.31\% \\
	0     & 0     & 1     & 1     & 2     & 56.25\% & 66.18\% \\
	0     & 0     & 1     & 1     & 3     & 56.20\% & 66.09\% \\
	0     & 1     & 0     & 0     & 1     & 70.85\% & 75.43\% \\
	0     & 1     & 0     & 0     & 2     & 68.19\% & 72.65\% \\
	0     & 1     & 0     & 0     & 3     & 68.35\% & 72.61\% \\
	0     & 1     & 0     & 1     & 1     & 57.68\% & 67.55\% \\
	0     & 1     & 0     & 1     & 2     & 55.74\% & 65.36\% \\
	0     & 1     & 0     & 1     & 3     & 55.88\% & 65.43\% \\
	0     & 1     & 1     & 0     & 1     & \textbf{71.91\%} & \textbf{76.37\%} \\
	0     & 1     & 1     & 0     & 2     & 69.44\% & 73.82\% \\
	0     & 1     & 1     & 0     & 3     & 69.42\% & 73.70\% \\
	0     & 1     & 1     & 1     & 1     & 59.09\% & 69.31\% \\
	0     & 1     & 1     & 1     & 2     & 57.49\% & 67.30\% \\
	0     & 1     & 1     & 1     & 3     & 57.29\% & 67.07\% \\
	1     & 0     & 0     & 0     & 1     & 36.40\% & 57.44\% \\
	1     & 0     & 0     & 0     & 2     & 40.78\% & 58.42\% \\
	1     & 0     & 0     & 0     & 3     & 40.85\% & 58.31\% \\
	1     & 0     & 0     & 1     & 1     & 40.88\% & 59.59\% \\
	1     & 0     & 0     & 1     & 2     & 45.99\% & 60.97\% \\
	1     & 0     & 0     & 1     & 3     & 45.68\% & 60.59\% \\
	1     & 0     & 1     & 0     & 1     & 34.45\% & 56.63\% \\
	1     & 0     & 1     & 0     & 2     & 39.73\% & 58.00\% \\
	1     & 0     & 1     & 0     & 3     & 39.73\% & 57.79\% \\
	1     & 0     & 1     & 1     & 1     & 38.56\% & 58.81\% \\
	1     & 0     & 1     & 1     & 2     & 43.64\% & 59.91\% \\
	1     & 0     & 1     & 1     & 3     & 43.97\% & 59.96\% \\
	1     & 1     & 0     & 0     & 1     & 38.92\% & 59.00\% \\
	1     & 1     & 0     & 0     & 2     & 42.52\% & 59.58\% \\
	1     & 1     & 0     & 0     & 3     & 42.72\% & 59.59\% \\
	1     & 1     & 0     & 1     & 1     & 43.99\% & 61.61\% \\
	1     & 1     & 0     & 1     & 2     & 47.16\% & 62.01\% \\
	1     & 1     & 0     & 1     & 3     & 47.29\% & 61.88\% \\
	1     & 1     & 1     & 0     & 1     & 37.17\% & 58.58\% \\
	1     & 1     & 1     & 0     & 2     & 41.38\% & 59.36\% \\
	1     & 1     & 1     & 0     & 3     & 41.70\% & 59.37\% \\
	1     & 1     & 1     & 1     & 1     & 41.78\% & 60.86\% \\
	1     & 1     & 1     & 1     & 2     & 45.68\% & 61.52\% \\
	1     & 1     & 1     & 1     & 3     & 46.10\% & 61.64\% \\
	\bottomrule
\end{longtable}%

De tabel leert ons dat een combinatie van een N-Gram tokenizer van maximum 1 woord, het stemmingsalgoritme en de stopwoordfiltering leidt tot het beste resultaat voor beide classificaties. Als we de frequentie weergeven in de woordvector, dan leidt dit meestal tot een daling van het aantal correct geclassificeerde documenten, terwijl de term frequency-inverse document frequency het aantal correct geclassificeerde documenten zelfs sterk doet dalen. 

\textbf{TODO: is hier een verklaring voor te vinden? Is dit afh van de dataset?}

Om het vector space model op te stellen voor de dataset van Het Laatste Nieuws, gaan we dus een N-Gram tokenizer van 1 woord gebruiken en zowel stemming als stopwoordfiltering toepassen. Term frequency-inverse document frequency en frequentie van woorden gebruiken we niet.

\subsection{Feature selection}\label{IG}
Het WEKA framework bevat een aantal mogelijkheden voor feature selection. Volgens \cite{Forman2003} zijn de 2 beste methodes voor feature selection \textit{Bi-Normal Seperation} (BNS) en \textit{Information Gain} (IG). 

Als $P_i$ de globale probabiliteit is van klasse $i$, en $p_{i}(w)$ de probabiliteit van klasse $i$ als het document het woord $w$ bevat. $F(w)$ is dan de fractie van de documenten die het woord $w$ bevatten. Information Gain $I(w)$ is als volgt gedefinieerd voor een gegeven woord  $w$:

\begin{math}
I(w) = -\sum_{i=1}^{k}P_i\cdot\log{(P_i)} + F(w)\cdot\log(p_i(w)) + (1 - F(w))\cdot\sum_{i=1}^{k}(1-p_i(w))\cdot\log(1-p_i(w))
\end{math}

Hoe groter de waardes voor IG $I(w)$, hoe groter de discriminerende kracht van het woord $w$. Voor een corpus met $n$ documenten en $d$ woorden is de complexiteit van IG $O(n \cdot d \cdot k)$.

We testen de implementatie van IG uit  in combinatie met het Ranker-algoritme in WEKA, op het woordenboek met de gekozen opties bij preprocessing. Dit levert de resultaten in tabel \ref{tab:feature-selection}.

\begin{table}[htbp]
	\centering
	\caption{Vergelijking voor en na feature selection met de Information Gain-methode. \textbf{TODO}}
	\begin{tabular}{rrrrrrr}
		\toprule
		\multicolumn{3}{c}{\textbf{Origineel}} & \multicolumn{1}{c|}{\textbf{}} & \multicolumn{3}{c}{\textbf{Information Gain}} \\
		\midrule
		\textbf{NB} & \textbf{NBM} & \textbf{woorden}    &   \multicolumn{1}{c|}{}   & \textbf{NB} & \textbf{NBM} & \textbf{woorden} \\
		71.91\% & 76.37\% & 1000  &   \multicolumn{1}{c|}{}    & 72.01\% & 76.51\% & 893 \\
		\bottomrule
	\end{tabular}%
	\label{tab:feature-selection}%
\end{table}%

We merken een sterke daling in het aantal woorden (10,70\%), maar tergelijkertijd een lichte stijging in het aantal juist geclassificeerde documenten (voor de "Naive Bayes$"$ classifier 0,10\%, voor de "Naive Bayes Multinomial" classifier 0.14\%).  We kunnen hier dus besluiten dat IG een algoritme is dat zeer goed werkt voor tekstuele data. De zogenaamde ruis wordt weggenomen uit de woordvector. Dit wordt dus ook toegepast na de preprocessing van de documentencollectie.


\subsection{Classificatie in hoofdcategorie\"en}
We starten met het classificeren in hoofdcategorie\"en. De dataset van januari 2008 bevat 11 hoofdcategorie\"en: Auto, Bizar, Geld, iHLN, Muziek, Nieuws, Reizen, Showbizz, Sport, Wetenschap en You. De verdeling van de artikels over deze hoofdcategorie\"en kan gevonden worden in tabel \ref{tab:hln-2008-01-cat}.

Om de best mogelijke classificatie te vinden, wordt een testbank opgesteld waarin een selectie van algoritmes voor classificatie in WEKA (46) opgenomen zijn. Deze classifiers worden met de standaardopties uitgevoerd op de data van januari 2008, die vooraf door de preprocessing-fase gegaan is. De classifiers worden vervolgens gerankschikt volgens het percentage correct geclassificeerde artikels, gebruik makend van een 10-fold cross-validation (zie \ref{10-fold-cv}). De resultaten zijn te vinden in \textbf{bijlage TODO}.

\subsubsection{Na\"ieve baseline classifier}\label{naieve-classifier}
Het aantal woorden die opgenomen wordt in de preprocessing van de documentencollectie wordt gevarieerd van 100 t.e.m. 20000 attributen met telkens een increment van 100 attributen. Nadien wordt de Information Gain methode (zie \ref{IG}) toegepast om eventuele ruis te verwijderen. Snel - na 2000 attributen - merken we echter dat enkele methodes het heel erg slecht doen. Een eerste zeer na\"ieve classifier waarmee vergeleken wordt is een classifier die alle artikels toekent aan de categorie met het meeste artikels (in dit geval nieuws). Deze na\"ieve implementatie zou in principe 31,69\% correct classificeren. In WEKA wordt deze ge\"implementeerd onder de naam ZeroR. Dit algoritme is ook opgenomen in lijst van classifiers. 

\begin{table}[htbp]
	\centering
	\caption{Classifiers die het slechter of even goed doen als de na\"ieve baseline classifier (ZeroR) bij classificatie van artikels in hoofdcategorie\"en na bijhouden van 2000 woorden in de woordvector.}
	\begin{tabular}{rrr}
		\toprule
		& 1000  & 2000 \\
		\midrule
		CVParameterSelection & \multicolumn{1}{c}{31.70\%} & \multicolumn{1}{c}{31.70\%} \\
		Grading & \multicolumn{1}{c}{31.70\%} & \multicolumn{1}{c}{31.70\%} \\
		MultiScheme & \multicolumn{1}{c}{31.70\%} & \multicolumn{1}{c}{31.70\%} \\
		Stacking & \multicolumn{1}{c}{31.70\%} & \multicolumn{1}{c}{31.70\%} \\
		StackingC & \multicolumn{1}{c}{31.70\%} & \multicolumn{1}{c}{31.70\%} \\
		Vote  & \multicolumn{1}{c}{31.70\%} & \multicolumn{1}{c}{31.70\%} \\
		ZeroR & \multicolumn{1}{c}{31.70\%} & \multicolumn{1}{c}{31.70\%} \\
		VFI   & \multicolumn{1}{c}{13.02\%} & \multicolumn{1}{c}{8.93\%} \\
		\bottomrule
	\end{tabular}%
	\label{tab:naive-classifiers}%
\end{table}

Als we ZeroR als baseline classifier beschouwen om te bepalen wat een slechte classifier is, dan vinden we enkele methodes die deze baseline niet halen en die het m.a.w. slechter of net even goed doen als ZeroR bij 2000 woorden. De zwakke classifiers zijn VFI (Voting Feature Intervals), Vote, Stacking, StackingC, MultiScheme, Grading, CVParameterSelection en ClassificationViaClustering. Als we enkele van deze algoritmes van dichterbij gaan bekijken, zien we dat o.a. Vote, VFI, Stacking, StackingC, MultiScheme en Grading enkele extra parameters verwachten waarin (enkele) andere classifiers worden opgegeven. Die methodes gaan dus andere classifiers combineren en worden verder besproken in \ref{combineren-classifiers}, waar de voordelen van het combineren van classifiers wordt onderzocht. 

CVParameterSelection is een classifier die een wrapper is rond een andere classifier. Deze functie is handig om te bepalen welke parameters van een bepaalde functie optimaal presteren. Omdat deze functie veel rekenkracht vraagt en op zich geen echte classificatie uitvoert, wordt deze hier niet verder behandeld. 

Het is voor het verder onderzoek naar de beste classifiers voor deze dataset niet nuttig dat de vermelde classifiers verder individueel onderzocht worden. Om rekenkracht te sparen worden deze dan ook niet meer opgenomen in de resultaten na 1000 woorden.

\subsubsection{One Rule}
Enkele methodes doen het net iets beter dan de na\"ieve ZeroR baseline, maar leveren geen verbetering op naar mate het aantal woorden toeneemt. We vergelijken de methodes nu met de \textit{OneR} methode. OneR is een simpel classificatie algoritme dat \'e\'en regel voor elk attribuut in de data genereert. Het selecteert dan de regel met de kleinste totale fout als zijn $"$one rule$"$. Om een regel op te stellen voor een attribuut, construereert OneR een frequentietabel voor elk attribuut tegenover het doelobject.

\begin{table}[htbp]
	\centering
	\caption{Classifiers die het slechter of even goed doen als de OneR methode bij classificatie van artikels in hoofdcategorie\"en na bijhouden van 2000 woorden in de woordvector.}
	\begin{tabular}{rrr}
		\toprule
		& \multicolumn{1}{c}{1000} & 2000 \\
		\midrule
		ConjunctiveRule & \multicolumn{1}{c}{49.64\%} & \multicolumn{1}{c}{49.60\%} \\
		AdaBoostM1 & \multicolumn{1}{c}{49.61\%} & \multicolumn{1}{c}{49.61\%} \\
		MultiBoostAB & \multicolumn{1}{c}{49.61\%} & \multicolumn{1}{c}{49.61\%} \\
		OneR  & \multicolumn{1}{c}{49.61\%} & \multicolumn{1}{c}{49.61\%} \\
		DecisionStump & \multicolumn{1}{c}{49.61\%} & \multicolumn{1}{c}{49.61\%} \\
		ClassificationViaClustering & 39.56\% & \multicolumn{1}{c}{39.38\%} \\
		\bottomrule
	\end{tabular}%
	
	\label{tab:addlabel}%
\end{table}%
 

De classifiers die het slechter of even goed doen als de OneR methode zijn: ConjunctiveRule, AdaBoostM1, MultiBoostAB, DecisionStump en ClassificationViaClustering. Opnieuw zijn enkele van deze classifiers wrappers rond anderen. AdaBoostM1 en MultiBoostAB zijn boosting algoritmes (zie \ref{boosting}) waarvan de default classifier DecisionStump is. Boosting kan een interessant gegeven zijn wanneer bepaalde klasses weinig gekozen worden door het classificatie algoritme.

\subsubsection{Naive Bayes}
Een volgend algoritme die als baseline gebruikt wordt is de Naive Bayes classifier bij 5000 woorden. Enkele methodes vallen af omdat ze te veel geheugen gebruiken en dus niet werken voor 5000 attributen. Dit zijn MulticlassClassifier, J48graft en SimpleCart. Hun maximaal behaalde percentage voor de documentencollectie, bij hun respectievelijk hoogst mogelijke aantal attributen kan terug gevonden worden in tabel \ref{tab:naive-bayes-out-of-memory}.

\textbf{TODO}

Enkele van de beschouwde methodes in deze stap zijn ook boosting- of combinatieclassifiers. LogitBoost en RacedIncrementalLogitBoost zijn beiden boosting classifiers (zie \ref{boosting}) en ClassificationViaRegression is een classifier die gebruikt wordt om bij stacking of grading (zie \ref{combineren-classifiers}) methodes de verschillende classifiers samen te voegen.

Resten nog de classifiers die het slechter of even goed doen dan de Naive Bayes methode bij 5000 woorden.  


\subsection{Combineren van classifiers voor hoofdcategorie\"en}
\label{combineren-classifiers}
\subsubsection{Voting}
Het is mogelijk dat we door verschillende goede individuele classifiers te combineren, een betere algemene classificatie bekomen. Een algoritme dat precies daarvoor gebruikt wordt is "Vote"  \cite{Kittler1998}. Dit is een algoritme dat zelfstandig geen classificatie uitvoert (zie ook resultaten individuele classifiers, \ref{naieve-classifier}). Het neemt een aantal zelf gekozen classifiers samen en laat die stemmen om te bepalen tot welke klasse een bepaalt document behoort. Om er zeker van te zijn of dit een effectieve verbetering oplevert passen we het "Vote$"$-algoritme toe met de top 4 beste algoritmes voor classificatie in hoofdcategorie\"en: DMNBText, MultinomialNaiveBayes, SMO en ComplementNaiveBayes. 

\begin{table}[htbp]
	\centering
	\caption{Aantal correct geclassificeerde artikels bij gebruik van het Vote-algoritme met classifiers DMNBText, MultinomialNaiveBayes, SMO en ComplementNaiveBayes voor dataset Het Laatste Nieuws (januari 2008) met verschillende combinatiemethodes.}
	\begin{tabular}{rr}
		\toprule
		Combinatiemethode & Correct geclassificeerd \\
		\midrule
		Average of Probabilities & 89,65\% \\
		Product of Probabilities & 86,45\% \\
		Majority Voting & 89,92\% \\
		Minimum Probability & 86,45\% \\
		Maximum Probability & 86,56\% \\
		\bottomrule
	\end{tabular}%
	\label{tab:vote}%
\end{table}%

Het Vote-algoritme laat toe om met verschillende methodes de originele classificaties te combineren. Deze worden voorgesteld in tabel \ref{tab:vote}, samen met het percentage correct geclassificeerde artikels voor de preprocessed dataset van januari 2008.

De methode "Majority Voting" geeft een licht beter resultaat dan de beste individuele classifier (DMNBText; 89,86\% correct geclassificeerd). Het verschil valt echter te verwaarlozen (0,06\%). 

\subsubsection{Stacking}\label{stacking}
Een andere methode om mogelijk een beter classificatie te bekomen is het "Stacking$"$-algoritme. Meestal presteert dit beter dan voting \cite{Sigletos2005}. \textit{Stacking} wordt gezien als een generalisatie van het "Vote$"$-algoritme. Het combineert de voorspellingen van verschillende andere classifiers en past een zelfgekozen combinatie-algoritme toe met de voorspellingen van de andere algoritmes als extra input. 

Het combineert meerdere classifiers die gegenereerd zijn door verschillende lerende algoritmes $L_1$,...,$L_n$ op een dataset S toe te passen. De dataset bestaat uit elementen in de vorm van $s_i = (x_i, y_i)$, welke in dit geval de paren van de feature vectoren ($x_i$) met hun classificatie ($y_i$) voorstellen. In de eerste fase worden base-level classifiers $C_1$,...,$C_n$ gegenereerd, waar $C_i = L_{i}(S)$. In de tweede fase wordt een meta-level classifier geleerd om de output van de base-level classifiers te combineren.


\textbf{TODO: resultaten}

\subsubsection{Grading}\label{grading}
Een andere methode om mogelijk een beter classificatie te bekomen is met het "Grading$"$-algoritme \cite{SEEWALD}. \textit{Grading} maakt gebruikt van \textit{graded} voorspellingen voor een trainingsset van meta classifiers die leren voorspellen wanneer de basisclassifier correct is. De trainingsset van deze meta classifiers is geconstrueerd gebruik makend van graded voorspellingen van de corresponderende basisclassifier als de nieuwe klasselabels voor de originele attributen. 

\textbf{TODO: resultaten}
\subsection{Ensemble learning voor subcategorie\"en}
\subsubsection{Boosting}
Boosting is een algemene methode om de performantie van lerende algoritmes te verbeteren. Het kan gebruikt worden om de fout van een zwak lerend algoritme dat consistent classifiers genereert die slechts enigszins beter zijn dan random gokken, significant te verlagen \cite{Freund1996}.


\textbf{TODO}

\subsubsection{Bagging}
\textbf{TODO}
 

\subsection{Classificatie in subcategorie\"en}
Het is mogelijk dat een classificatie in subcategorie\"en een ander algoritme vereist dan het classificeren in hoofdcategorie\"en. Daarom testen we ook voor de subcategorie\"en van hoofdcategorie "Nieuws" welk algoritme het meeste documenten kan classificeren. Per hoofdcategorie zijn heel wat subcategorie\"en opgenomen, de verdeling van de artikels over de subcategorie\"en wordt voorgesteld in tabel \ref{tab:hln-2008-01-nieuws}. Voor het testen van de classifiers is hier ook telkens de dataset van januari 2008 gebruikt. 

\begin{table}[htbp]
	\centering
	\caption{Add caption}
	\begin{tabular}{rrr}
		\toprule
		\textbf{Subcategorie} & \textbf{Aantal} & \textbf{Gewicht} \\
		\midrule
		Binnenland & 1582  & 49.28\% \\
		Buitenland & 1533  & 47.76\% \\
		Israel-Palestina & 34    & 1.06\% \\
		Verdwijning Maddie & 24    & 0.75\% \\
		Onderwijs & 17    & 0.53\% \\
		Islam & 8     & 0.25\% \\
		Ronald Janssen & 6     & 0.19\% \\
		Irak  & 4     & 0.12\% \\
		Hoe moet ik solliciteren? & 1     & 0.03\% \\
		Zaak Holloway & 1     & 0.03\% \\
		\textbf{Totaal} & \textbf{3210} & \textbf{100.00\%} \\
		\bottomrule
	\end{tabular}
	\label{tab:hln-2008-01-nieuws}%
\end{table}%

\subsection{Combineren van classifiers voor subcategorie\"en}
\subsubsection{Voting}
Om tot een betere classificatie te komen voor subcategorie\"en, worden de beste 8 classifiers gebruikt bij 9900 woorden: DMNBText, Dagging,  RandomCommittee, SMO, Logistic, RandomSubSpace, NaiveBayesMultinomial en ComplementNaiveBayes. De verschillende combinatiemethodes voor het Vote-algoritme werden getest en kunnen gevonden worden in tabel \ref{tab:vote-sub}. 



\subsubsection{Stacking}
\subsubsection{Grading}
\subsection{Ensemble learning voor subcategorie\"en}
\subsubsection{Bagging}
\subsubsection{Boosting}

\subsection{Relatie tussen classificatie en grootte van de dataset}
\textbf{TODO}

\section{Clustering}
Voor clustering wordt gebruikt gemaakt van het "K-means$"$ algoritme en het "Expectation Maximization" (EM) algoritme. 

\section{Aanbevelingssysteem}
\chapter{Resultaten van het aanbevelingssysteem}\label{hs:resultaten}
\section{Baseline recommender system}
Om de resultaten van het nieuw ontworpen systeem te vergelijken met bestaande content-based recommender systemen, wordt gebruik gemaakt van een baseline recommender gebasseerd op term frequency-inverse document frequency. Deze baseline recommender werkt bovenop Lenskit.

\subsection{Illustratie van de werking van een baseline recommender}

\section{Vergelijking van het baseline recommender systeem en de nieuw ontworpen recommender}

% appendices
\appendix

% hier worden de appendices ingevoegd (\includes)


\bibliography{bibliografie/ugent-bjorn-2016}{}
\bibliographystyle{plain}

\backmatter

% eventueel: lijst van figuren en tabellen
\listoffigures
\listoftables

% lege pagina (!!)


% kaft

\end{document}
\grid
